\scanpage{272}
\mspage{13}

3. \deleted{\uline{\unsure{Composition} des intervalles orientés}} \uline{Chemins}. \uline{Catégorie des chemins d'un type donné}.

Soient $I$ un intervalle orienté, et $X$ une catégorie. On appelle \uline{chemin de type} $I$ dans $X$ \unsure{tout} foncteur $I\unsure{\longrightarrow|\text{---}}X$, \deleted{\unsure{et le}} et \uline{catégorie des chemins de type} $I$ \uline{dans} $X$ la catégorie

\[\underline{\mathrm{Ch}}_I(X)=\underline{\mathrm{Hom}}(I,X).\]

Comme $I$ est \uline{canoniquement} isom. à une catégorie intervalle standard $I_0$, \deleted{le ``type''} \deleted{\unsure{les données}} on peut identifier les chemins de type $I$, ou de type $I_0$. Le type \inserted{$T$} d'un chemin \inserted{de longueur $n$ ($=\operatorname{card}I$)} dans $X$ peut donc être décrit par une suite de $n$ signes $+$ ou $-$

\[\unsure{\tau|T|\xi|\varepsilon}_1,\tau_2,\ldots,\tau_n\qquad(\tau_i\in\{-1,+1\}\text{ pour }1\leq i\leq n).\]

À cause du fait que $I$ est engendrée par ses flèches de transition, un foncteur $c:I\unsure{-|\longrightarrow}X$ est dét\inserted{\unsure{er}}miné \deleted{de façon univ} quand on connaît

a) les $c(i)$ ($i\in I$) i.e une suite de \inserted{$n+1$} objets de $X$ :

\[\deleted{\unsure{a|x}}A_0,\deleted{a}A_1,\ldots,\deleted{a}A_n\,;\]

% slanted
\margin{$\unsure{A|a}_0$ est dit la \uline{source} \unsure{du chemin|chemin}, $a_n$ est son \uline{but}.}

b) Les \deleted{\unsure{fl}} $c(\underline{u}_i)$ (où les $\underline{u}_i$ sont les flèches de transition de $I$) i.e. pour tout $i\unsure{\notin|\in}[1,n]$ une flèche \inserted{$u_i$} dans $X$ entre $\deleted{a}A_{i-1}$ et $\deleted{a}A_i$, allant de $\deleted{a}A_{i-1}$ dans $\deleted{a}A_i$ si $\underline{u}_i$ est directe ($\tau_i=+1$), et de $\deleted{a}A_i$ dans $\deleted{a}A_{i-1}$ si $\underline{u}_i$ est rétrograde ($\tau_i=-1$),

et \deleted{aux} \inserted{à toute telle} donnée a), b) correspond un foncteur. Ainsi on pourra identifier les chemins de type $\deleted{\unsure{I}}\unsure{T|I}$

% archive sheet number: 271
